\section{Trabajo futuro}
\label{sec:future-work}

\begin{itemize}

\item \textbf{Familias adicionales validadas.} La extensi\'on de mayor
  valor (\S\ref{sec:limitations}, punto~\ref{lim:n3}) es una validaci\'on
  completa respaldada por ingenier\'ia inversa manual de m\'as
  familias. Ya se realiz\'o una comprobaci\'on m\'as ligera e informal
  frente a una 4\textsuperscript{a} familia sin afinar (AsyncRAT): una
  detonaci\'on real en CAPE m\'as el reenv\'io autom\'atico de sus
  cargas soltadas confirm\'o que la pipeline produce un resultado
  sensato y sin ca\'idas de principio a fin. Esto carece de datos de
  referencia manuales y no se trata como evidencia de validaci\'on, solo
  como una comprobaci\'on interna de cordura de que la pipeline
  generaliza en absoluto (documentada en \code{README.md}); no
  sustituye a la extensi\'on completa con datos de referencia que
  describe este punto.

\item \textbf{Ampliar a\'un m\'as la cobertura de reglas est\'aticas.}
  Las 86 de 474 t\'ecnicas relevantes para Windows de
  \S\ref{sec:coverage} dejan aproximadamente 388 sin cubrir. Cerrar esa
  brecha es un trabajo deliberadamente lento, t\'ecnica a t\'ecnica, no
  un conjunto de tareas para hacer en bloque: cada adici\'on necesita
  su propia justificaci\'on a nivel de descripci\'on de MITRE y una
  nueva ejecuci\'on completa frente a las muestras validadas para
  descartar nuevos falsos positivos, el mismo list\'on que se aplic\'o
  a la ampliaci\'on original de 30 a 86.

\item \textbf{An\'alisis din\'amico Linux/ELF para familias de botnets
  IoT.} Investigado de forma concreta y no descartado por suposici\'on,
  y aplazado deliberadamente por tres razones espec\'ificas, cada una
  verificada de forma independiente: el analizador ELF actual de la
  pipeline extrae solo metadatos estructurales b\'asicos, sin ninguna
  regla de detecci\'on ATT\&CK conectada, a diferencia de la v\'ia PE
  con sus aproximadamente 86 t\'ecnicas mapeadas; la propia
  documentaci\'on de CAPE describe su soporte de sandboxing Linux/ELF
  como una capacidad en desarrollo activo, no de nivel producci\'on; y
  todav\'ia no existe ninguna muestra de familia ELF con datos de
  referencia frente a la que verificar cualquier nueva regla de
  detecci\'on espec\'ifica de Linux, sin lo cual construir reglas
  repetir\'ia exactamente el riesgo de sobreajuste que ya se\~nala
  \S\ref{sec:limitations} para el lado Windows. Dado lo
  desproporcionadamente comunes que son las botnets basadas en ELF en
  dispositivos IoT y Linux embebido en comparaci\'on con la atenci\'on
  que les presta la propia literatura de investigaci\'on de malware,
  se considera una extensi\'on genuinamente valiosa, pero una que
  primero necesita una muestra de familia ELF validada y una capacidad
  real de an\'alisis din\'amico, no un conjunto de reglas reajustado
  sobre el an\'alisis est\'atico por s\'i solo.

\item \textbf{Un estudio de calibraci\'on de confianza con potencia
  estad\'istica.} El hallazgo de \S\ref{sec:confidence-calibration} de
  que 16 de las 17 discrepancias auditadas provienen de hallazgos de
  una sola fuente, con la \'unica excepci\'on entre fuentes
  (\technique{T1497} de RoningLoader) atribuible a una cuesti\'on de
  completitud de los datos de referencia y no a un nivel de confianza
  mal calibrado, es una se\~nal direccional de tres muestras, no una
  afirmaci\'on con potencia estad\'istica, y se beneficiar\'ia de un
  conjunto de validaci\'on m\'as grande y dise\~nado para ello.

\item \textbf{Recuperaci\'on de claves en tiempo de ejecuci\'on
  mediante emulaci\'on.} La clave RC4 que nunca existe como bytes
  contiguos en la DLL de RoningLoader (\S\ref{sec:limitations},
  punto~\ref{lim:key-ceiling}) es, en principio, recuperable emulando su rutina de
  ensamblado (p.\ ej.\ con Unicorn) en lugar de mediante cualquier
  b\'usqueda de patr\'on de bytes est\'atica, un esfuerzo de
  ingenier\'ia sustancialmente mayor que el m\'odulo actual de
  recuperaci\'on est\'atica de configuraci\'on, y fuera de alcance de
  esta tesis en consecuencia.

\end{itemize}
